\documentclass[pdflatex,sn-mathphys-num]{sn-jnl}

\usepackage{amsmath,amssymb,amsfonts}%
\usepackage{amsthm}%
\usepackage{enumitem}%

%% theorem styles follow sn-jnl presets
\theoremstyle{thmstyleone}%
\newtheorem{theorem}{Theorem}[section]
\newtheorem{proposition}[theorem]{Proposition}
\newtheorem{lemma}[theorem]{Lemma}
\newtheorem{corollary}[theorem]{Corollary}

\theoremstyle{thmstylethree}%
\newtheorem{definition}[theorem]{Definition}

\theoremstyle{thmstyletwo}%
\newtheorem{remark}[theorem]{Remark}
\newtheorem{example}[theorem]{Example}
\newtheorem{openproblem}{Open Problem}

%% Math operators
\DeclareMathOperator{\Sym}{Sym}
\DeclareMathOperator{\tr}{tr}
\DeclareMathOperator{\Hess}{Hess}
\DeclareMathOperator{\image}{im}
\DeclareMathOperator{\coker}{coker}
\DeclareMathOperator{\Diff}{Diff}
\DeclareMathOperator{\GL}{GL}
\DeclareMathOperator{\UU}{U}
\DeclareMathOperator{\OO}{O}
\DeclareMathOperator{\Met}{Met}
\DeclareMathOperator{\Herm}{Herm}
\DeclareMathOperator{\id}{id}
\DeclareMathOperator{\Ric}{Ric}
\DeclareMathOperator{\Rm}{Rm}
\DeclareMathOperator{\Hol}{Hol}
\DeclareMathOperator{\SO}{SO}
\DeclareMathOperator{\SU}{SU}
\DeclareMathOperator{\diam}{diam}
\DeclareMathOperator{\Hom}{Hom}
\DeclareMathOperator{\Vol}{Vol}
\DeclareMathOperator{\Crit}{Crit}
\DeclareMathOperator{\Lev}{Lev}

%% Complex analysis shorthands
\newcommand{\dbar}{\bar{\partial}}
\newcommand{\pphi}{\partial\varphi}
\newcommand{\dbarphibar}{\dbar\bar\varphi}
\newcommand{\Gvarphi}{i\,\partial\varphi \wedge \dbar\bar\varphi}
\newcommand{\Levphi}[1]{i\partial\dbar\,#1(\varphi)}
\newcommand{\omegah}{\omega_h}
\newcommand{\Jbundle}{\mathcal{J}}

\begin{document}

\title[Complex Seam Geometry]{Complex Seam Geometry: A Jet-Theoretic Framework
for the Hermitian Generation of Geometric Structures}

\author*[1]{\fnm{Lars} \sur{R\"onnb\"ack}}\email{lars@uptochange.com}

\affil*[1]{\orgdiv{Department of Computer and Systems Sciences},
  \orgname{Stockholm University},
  \orgaddress{\city{Stockholm}, \country{Sweden}}}

\abstract{%
We develop a jet-theoretic framework for the seam-rule paradigm in
the setting of \emph{complex-valued scalar fields} on Hermitian manifolds.
Viewing a rule as a $\UU(m)$-equivariant natural differential operator on
complex jet bundles yields a classification of all isotropic
Hermitian-metric-generating rules of order $\le 2$: an explicit
\emph{seven-real-parameter family} consisting of a conformal term,
a unitary-isotropic quadratic gradient family, and a complex Levi-form
term, with coefficients depending smoothly on $\UU(m)$-invariant scalar
jets.  The Levi-form generator $\operatorname{Re}(\gamma\,i\partial\dbar\varphi)$
is genuinely new relative to the real-seam framework and arises from the
tighter $\UU(m)$-equivariance replacing $\OO(2m)$.

We then develop the deformation theory of complex rules.  The
linearisation at a seam is a real-linear differential operator whose
ellipticity character is governed by the Levi-form term; the
degeneration locus unifies the holomorphic critical-point locus and the
Levi-degenerate set of the seam.

Finally, we study a universality problem: when can complex scalar seam
data generate an open neighbourhood of Hermitian metrics up to
biholomorphism?  We prove that the Levi-form generator achieves local
universality within a fixed K\"ahler class via the $\partial\dbar$-lemma,
and state open problems for the full complex rule.
}

\keywords{complex seam geometry, natural differential operators,
Hermitian metrics, Levi form, K\"ahler geometry,
$\UU(m)$-equivariance}

\pacs[MSC Classification]{53C55 (primary), 58A20, 53C21, 53A35,
32Q15, 53C25, 32W05}

\maketitle

% ================================================================
\section{Introduction}
% ================================================================

The observation that drives this paper is elementary but has
consequences that reach into several applied domains.  In the
real-seam framework of~\cite{Ronnback2025}, the naturality group is
$\OO(n)$ (the orthogonal group of the background Riemannian metric),
and $\OO(n)$-invariant theory yields exactly three generators of
isotropic metric rules at order $\le 2$.  When the ambient manifold
carries a compatible complex structure---making it a Hermitian
manifold $(M^{2m}, h, J)$---and the seam is promoted to a
\emph{complex}-valued function $\varphi \in C^\infty(M,\mathbb{C})$,
two things happen simultaneously:

\begin{enumerate}[label=(\roman*)]
\item The naturality group narrows from $\OO(2m)$ to $\UU(m)$
  (the unitary group), because coordinate changes must now preserve
  both $h$ and $J$.  A tighter equivariance constraint means
  \emph{more} equivariant maps are admissible, not fewer: the
  finer orbit structure of $\UU(m)$ distinguishes holomorphic
  from antiholomorphic data, allowing new generators that are
  invisible to $\OO(2m)$.
\item The jet of $\varphi$ decomposes under $\UU(m)$ into
  \emph{bidegree components} $(\partial\varphi, \dbar\varphi,
  \partial^2\varphi, \partial\dbar\varphi, \dbar^2\varphi)$
  rather than the single symmetric Hessian of the real case.
  The mixed component $\partial\dbar\varphi$ is the
  \emph{Levi form} of $\varphi$, and it generates a
  fundamentally new class of Hermitian metrics with no real
  analogue.
\end{enumerate}

These two points together produce the \emph{complex Levi-form generator}
in the classification theorem (Theorem~\ref{thm:classification-complex}).  The
Levi-form generator is why K\"ahler geometry---in which every
metric arises locally as $i\partial\dbar\psi$ for a real potential
$\psi = \mathrm{Re}(\varphi)$---is richer and more rigid than the
purely conformal or purely gradient-tensor approaches: it exploits
a generator that is native to the complex structure.

\subsection*{The recurrent construction in complex settings}

The same schema as in the real case appears pervasively in contexts
where the natural data is complex-valued:

\begin{center}
\emph{complex scalar field $\varphi$}
$\;\xrightarrow{\;\text{rule}\;}\;$
\emph{Hermitian structure}.
\end{center}

Concrete instances include:

\begin{itemize}
\item \textbf{K\"ahler metrics}: $\omega = i\partial\dbar\psi$
  for a strictly plurisubharmonic real potential $\psi$---the
  Levi-form generator with $\gamma = 1$ and all gradient/conformal
  coefficients set to $0$.
\item \textbf{Fubini--Study metric on $\mathbb{CP}^m$}:
  $\omega_{\mathrm{FS}} = i\partial\dbar\log(1 + |z|^2)$---a
  Levi-form rule applied to the log-norm seam.
\item \textbf{Fisher--Rao metric of a complex exponential
  family}: the log-partition function of an exponential family
  parametrised by complex $\theta$ is the seam; its Levi form
  $i\partial\dbar A(\theta)$ gives the Fisher information metric on the
  complex parameter manifold.
\item \textbf{Phase geometry in MRI and radar}: the k-space
  phase field $\varphi = \rho e^{i\phi}$ is the natural complex
  seam; the gradient-tensor generator $i\partial\varphi \wedge \dbar\bar\varphi$
  encodes the local frequency content as a Hermitian form.
\item \textbf{Quantum state geometry}: for a (local section of a)
  wave function $\psi \in L^2$, the Fubini--Study metric arises
  from the seam $\varphi = \log\psi$ via the Levi-form rule.
  The degeneration locus is the zero set of $\psi$---the nodal
  hypersurface---which plays the role of the critical-point locus
  in the real-seam framework.
\end{itemize}

\subsection*{Three payoffs of the jet-theoretic viewpoint}

As in the real case, the naturality framework delivers three
structural payoffs:

\begin{enumerate}[label=(\roman*)]
\item \textbf{Classification.}  $\UU(m)$-equivariant invariant theory
  constrains the space of possible rules to a seven-real-parameter family
  at order $\le 2$ under the assumption of at most linear dependence on the
  mixed Hessian.  Relative to the real-seam framework, there are two distinct
  sources of enrichment: (a) the mixed Hessian component $\partial\dbar\varphi$
  contributes the genuinely complex Levi-form generator, and (b) the 1-jet of a
  complex seam carries two independent holomorphic gradients
  $(\partial\varphi,\partial\bar\varphi)$, yielding a full $2\times 2$ Hermitian
  family of unitary-isotropic quadratic gradient terms
  (Theorem~\ref{thm:classification-complex}).
\item \textbf{Linearisation and inverse design.}  The linearised rule
  at a complex seam has a principal symbol involving $\xi \otimes \bar\xi$
  (from the Levi-form generator) rather than the symmetric $\xi \otimes \xi$
  of the real Hessian generator.  This changes the ellipticity character:
  the relevant operator is the \emph{complex Laplacian} $\partial\dbar$
  rather than the real Laplacian $\nabla^2$
  (\S\ref{sec:linearisation}).
\item \textbf{Universality.}  The image of the complex seam map
  is constrained to Hermitian metrics lying in the image of
  $i\partial\dbar$ (for the Levi-form rule) and to specific subsets of
  the K\"ahler cone.  The universality question---when can a complex seam
  generate an open neighbourhood of Hermitian metrics up to
  biholomorphism---has a richer answer than the real case, connected
  to the Calabi--Yau conjecture and $\partial\dbar$-lemma techniques
  (\S\ref{sec:universality}).
\end{enumerate}

\begin{remark}[Relation to the real-seam framework]
\label{rem:real-embedding}
The real-seam framework of~\cite{Ronnback2025} embeds into the
present one by restricting to real-valued seams $\varphi = s \in
C^\infty(M, \mathbb{R})$.  In this case the imaginary Levi-form
generator $i\partial\dbar\,\mathrm{Im}(\varphi)$ vanishes, and
$i\partial\dbar\,\mathrm{Re}(\varphi) = i\partial\dbar s$ coincides with
(a scalar multiple of) the $(1,1)$-part of the Riemannian Hessian
$\nabla^2_h s$ on a K\"ahler manifold.  The gradient-tensor generator
$i\partial s \wedge \dbar s$ is proportional to the $(1,1)$-form
representation of $ds \otimes ds$.  Thus the five complex generators
specialise to the three real generators when the seam is real, with
the Levi form and Hessian generators identified through the K\"ahler
structure.



In holomorphic coordinates,
\[
i\partial s \wedge \dbar s
= \frac{i}{2}\,\partial_\alpha s\,\partial_{\bar\beta} s\,dz^\alpha\wedge d\bar z^\beta,
\]
so the associated Hermitian matrix is the rank-$1$ outer product
$(\partial_\alpha s)(\partial_{\bar\beta} s)$, i.e. the $(1,1)$-projection
of $ds\otimes ds$.
\end{remark}

\begin{theorem}[Reduction to the real-seam framework]
\label{thm:reduction-real}
If the seam $\varphi$ is real-valued, $\varphi = s \in C^\infty(M,\mathbb{R})$,
then the complex universal decomposition~\eqref{eq:universal-complex} reduces to
the three-generator decomposition of the real-seam framework of~\cite{Ronnback2025}.
\end{theorem}

\begin{proof}[Proof sketch]
If $\varphi=s$ is real, then $\bar\varphi=\varphi$ and the two holomorphic
gradients coincide: $\partial\bar\varphi=\partial\varphi$.
Consequently the unitary-isotropic gradient family collapses to the single
generator $i\partial s\wedge\dbar s$.
Moreover $\mathrm{Im}(\varphi)=0$, so the imaginary Levi form vanishes and the
Levi-form generator reduces to $i\partial\dbar s$, which on a K\"ahler manifold
is (up to a scalar multiple) the $(1,1)$-part of the real Hessian generator.
These identifications are explained in Remark~\ref{rem:real-embedding}.
\end{proof}

% ================================================================
\section{Complex Seams and Their Jets}
\label{sec:seams}
% ================================================================

Throughout, $(M^{2m}, h, J)$ denotes a smooth Hermitian manifold of
complex dimension $m \ge 2$ (real dimension $n = 2m$), where $h$ is
a Riemannian metric and $J$ is a compatible complex structure.
We write $\omegah$ for the fundamental $(1,1)$-form of the Hermitian
structure:
\[
\omegah(X,Y) = h(JX, Y),
\quad X, Y \in TM.
\]
We use holomorphic local coordinates $(z^1, \ldots, z^m)$ throughout,
and adopt the standard bidegree convention $\partial = \partial_z$,
$\dbar = \partial_{\bar z}$.

\begin{definition}[Complex seam]
\label{def:seam-complex}
A \emph{complex seam} on $(M, h, J)$ is a smooth complex-valued
function $\varphi \in \mathcal{S}_\mathbb{C}(M) \subseteq
C^\infty(M, \mathbb{C})$, where $\mathcal{S}_\mathbb{C}(M)$ is an
admissible function class appropriate to the application (e.g.\
all smooth, plurisubharmonic, holomorphic, or discrete vertex-based
functions).  We write $\varphi = u + iv$ with $u,v \in C^\infty(M,\mathbb{R})$.
\end{definition}

\begin{remark}[The role of the decomposition $\varphi = u + iv$]
The real and imaginary parts $u = \mathrm{Re}(\varphi)$ and
$v = \mathrm{Im}(\varphi)$ are not symmetric in the complex
structure: under a holomorphic change of local coordinates,
$\partial u$ and $\partial v$ are related by the Cauchy--Riemann
operator, which encodes the complex structure $J$.  When $\varphi$
is holomorphic ($\dbar\varphi = 0$), the Cauchy--Riemann equations
force $v$ to be the harmonic conjugate of $u$, and the two Levi
forms $i\partial\dbar u$ and $i\partial\dbar v$ become related
(both vanish for a holomorphic seam, since $\partial\dbar\varphi = 0$).
This coupling between $u$ and $v$ is precisely what the complex
naturality group $\UU(m)$---and not merely $\OO(2m)$---detects.
\end{remark}

The Peetre--Slov\'ak localisation theorem~\cite{Peetre1960,Slovak1988}
applies verbatim: every smooth natural operator of finite order
factors through the complex jet bundle.

\begin{definition}[Complex $k$-jet of a seam]
\label{def:jet-complex}
Let $\varphi \in C^\infty(M, \mathbb{C})$.  The \emph{complex
$k$-jet of $\varphi$ at $x \in M$}, denoted $j^k_{x,\mathbb{C}}\varphi$,
is the equivalence class of $\varphi$ under the relation of having the
same Taylor expansion in holomorphic coordinates at $x$ through
bidegree order $(p,q)$ with $p + q \le k$.  In holomorphic
local coordinates $(z^\alpha, \bar{z}^{\bar\alpha})$,
\[
j^2_{x,\mathbb{C}}\varphi
= \Bigl(
  \varphi(x),\;\;
  \partial_\alpha\varphi(x),\;\;
  \dbar_{\bar\alpha}\varphi(x),\;\;
  \partial_\alpha\partial_\beta\varphi(x),\;\;
  \partial_\alpha\dbar_{\bar\beta}\varphi(x),\;\;
  \dbar_{\bar\alpha}\dbar_{\bar\beta}\varphi(x)
  \Bigr).
\]
The collection of all complex $k$-jets forms the \emph{complex
$k$-jet bundle} $J^k_\mathbb{C}(M, \mathbb{C}) \to M$.
\end{definition}

\begin{remark}[Bidegree decomposition of the 2-jet]
\label{rem:bidegree-decomp}
At order $k = 2$, the complex jet $j^2_{x,\mathbb{C}}\varphi$
decomposes under the $\UU(m)$-action (evaluated in $h$-unitary
frames at $x$) into irreducible components as follows.
Write $\varphi_0 = \varphi(x) \in \mathbb{C}$,
$v_\alpha = \partial_\alpha\varphi(x) \in T^{*1,0}_xM \cong \mathbb{C}^m$,
$\bar w_{\bar\alpha} = \dbar_{\bar\alpha}\varphi(x) \in T^{*0,1}_xM
\cong \overline{\mathbb{C}^m}$,
$P_{\alpha\beta} = \partial_\alpha\partial_\beta\varphi(x) \in
\Sym^2(T^{*1,0}_xM)$,
$B_{\alpha\bar\beta} = \partial_\alpha\dbar_{\bar\beta}\varphi(x) \in
T^{*1,0}_xM \otimes T^{*0,1}_xM$,
and
$\bar Q_{\bar\alpha\bar\beta} = \dbar_{\bar\alpha}\dbar_{\bar\beta}\varphi(x)
\in \Sym^2(T^{*0,1}_xM)$.
The irreducible $\UU(m)$-components are:
\begin{enumerate}[label=(\alph*)]
\item \textbf{Value:} $\varphi_0 \in \mathbb{C}$
  (trivial complex representation, $\mathrm{U}(1)$-charged).
\item \textbf{Holomorphic gradient:} $v = (v_\alpha) \in
  T^{*1,0}_xM$ (standard complex representation of $\UU(m)$).
\item \textbf{Antiholomorphic gradient:} $\bar w = (\bar w_{\bar\alpha}) =
  \dbar\varphi(x) \in T^{*0,1}_xM$ (conjugate representation).
\item \textbf{Holomorphic Hessian:} $P \in \Sym^2(T^{*1,0}_xM)$
  (symmetric square of the standard representation).
\item \textbf{Mixed (Levi) Hessian:} $B \in
  T^{*1,0}_xM \otimes T^{*0,1}_xM$ (tensor product of standard and
  conjugate).  This is the $\UU(m)$-representation containing the
  Hermitian forms as a real subspace; it further decomposes as
  $\mathbb{R}\,\omega_h \oplus \mathfrak{u}(m)^\perp$
  (trace plus trace-free Hermitian part).
\item \textbf{Antiholomorphic Hessian:} $\bar Q \in
  \Sym^2(T^{*0,1}_xM) = \overline{\Sym^2(T^{*1,0}_xM)}$
  (conjugate of the holomorphic Hessian representation).
\end{enumerate}
This six-component decomposition is substantially richer than the
three-component real decomposition
$j^2_x s = (s_0, \nabla s, \nabla^2 s)$, reflecting the finer
orbit structure of $\UU(m)$ compared with $\OO(2m)$.  The
components $v$ and $\bar w$ are \emph{independent} for a general complex
seam.  Indeed, $\overline{v} = \dbar\,\bar\varphi$ whereas
$\overline{\bar w} = \partial\,\bar\varphi$, so $\bar w$ coincides with
$\overline{v}$ only when $\varphi$ is real-valued.
In particular, the 1-jet carries two independent holomorphic gradients
$\partial\varphi$ and $\partial\bar\varphi$ (equivalently
$\partial\bar\varphi = \overline{\dbar\varphi}$), which yields a richer
family of unitary-isotropic quadratic gradient terms than in the real-seam
case.  Similarly $\bar Q = \overline{P}$.
The genuinely new data relative to the real case is the
decomposition of the Hessian into the holomorphic part~$P$,
the antiholomorphic part~$\bar Q$, and the mixed part~$B$.
\end{remark}

\begin{remark}[Comparison with the real 2-jet]
\label{rem:comparison}
Viewing $(M,h,J)$ as a real Riemannian manifold $(M^{2m},h)$
with $\OO(2m)$-symmetry, the real 2-jet $j^2_x\varphi$ (treating
$\varphi$ as a pair of real functions $(u,v)$) has three components:
value, real gradient, and real Hessian.  The $\UU(m)$ bidegree
decomposition of Remark~\ref{rem:bidegree-decomp} is a refinement:
it splits the real gradient into holomorphic and antiholomorphic
parts, and the real Hessian into holomorphic, mixed (Levi), and
antiholomorphic parts.  The Levi component $B$ is the critical
new actor: it is the part of the Hessian most directly controlled
by the complex structure, and it generates the $i\partial\dbar$
operator that does not appear in the real-seam classification.
\end{remark}

% ================================================================
\section{Rules as Natural Differential Operators}
\label{sec:rules}
% ================================================================

The natural target bundle for a complex rule is the bundle of real
$(1,1)$-forms:
\[
\mathcal{H}(M) := \Lambda^{1,1}(M) \cap \Lambda^2(M,\mathbb{R}).
\]
Inside it sits the open convex cone of positive Hermitian forms:
\[
\mathcal{H}^+(M) \;:=\;
\bigl\{\,\Omega \in \Lambda^{1,1}(M) \cap \Lambda^2(M,\mathbb{R}) \;:\;
\Omega \text{ is positive definite}\bigr\}.
\]
Thus $\mathcal{H}^+(M) \subseteq \mathcal{H}(M)$.
Locally, $\Omega \in \mathcal{H}^+(M)$ is a real $(1,1)$-form with
positive-definite coefficient matrix $h_{\alpha\bar\beta}(x)$
satisfying $\Omega = \frac{i}{2}h_{\alpha\bar\beta}\,dz^\alpha \wedge
d\bar{z}^\beta$ and $h_{\alpha\bar\beta} = \overline{h_{\beta\bar\alpha}}$,
$h > 0$.

\begin{definition}[Complex rule of order $k$]
\label{def:rule-complex}
A \emph{complex rule of order $k$} is a smooth bundle map
\begin{equation}
\label{eq:rule-complex}
\Phi \colon J^k_\mathbb{C}(M, \mathbb{C}) \times \mathcal{B}(M)
\longrightarrow \mathcal{H}(M)
\end{equation}
satisfying:
\begin{enumerate}[label=(R\arabic*)]
\item \label{item:locality-c}
  \textbf{Locality.}
  $\Phi$ depends fibrewise on $(j^k_{x,\mathbb{C}}\varphi,\, b_x)$
  where $b_x$ is the fibre of the background structure
  $\mathcal{B}(M)$ at $x$.
\item \label{item:naturality-c}
  \textbf{Naturality.}
  For every biholomorphism $\psi \colon M \to M$ preserving
  $\mathcal{B}$,
  \[
  \Phi\bigl(j^k_\mathbb{C}(\psi^*\varphi),\, \psi^*b\bigr)
  = \psi^*\,\Phi(j^k_\mathbb{C}\varphi,\, b).
  \]
  If $\mathcal{B}$ includes a Hermitian structure $(h,J)$, naturality
  reduces to $\UU(m)$-equivariance at each fibre.
\item \label{item:regularity-c}
  \textbf{Regularity.}
  $\Phi$ is smooth as a map of fibre variables.
\end{enumerate}
\end{definition}

\begin{definition}[Metric-admissible rule]
\label{def:admissible-rule}
A complex rule $\Phi$ is \emph{metric-admissible} on an open subset
$\mathcal{U} \subseteq J^k_\mathbb{C}(M, \mathbb{C}) \times \mathcal{B}(M)$ if
$\Phi(\mathcal{U}) \subseteq \mathcal{H}^+(M)$.  When the domain of interest is
clear from context, we simply say that $\Phi$ is a (Hermitian) metric rule.
\end{definition}

\begin{remark}[Why $\UU(m)$, not $\OO(2m)$]
\label{rem:unitary}
The naturality condition in axiom~\ref{item:naturality-c} is stated
for biholomorphisms, which preserve both $h$ and $J$.  On a
Hermitian manifold $(M^{2m},h,J)$, the stabiliser of the background
structure at a fibre is the unitary group $\UU(m)$, not the full
orthogonal group $\OO(2m)$.  This is the decisive difference from
the real-seam framework: $\OO(2m)$-equivariance would allow rules
that mix the $(1,0)$ and $(0,1)$ components of the jet, producing
no new structure beyond what is already visible in real coordinates.
$\UU(m)$-equivariance forbids such mixing, enforcing a bidegree
conservation that is ultimately responsible for the extra generator
in the classification theorem.

Concretely: the real Hessian $\nabla^2_h\varphi$ is a section of
$\Sym^2(T^*M \otimes \mathbb{C})$, which decomposes under $\UU(m)$
as $\Sym^2(T^{*1,0}M) \oplus (T^{*1,0}M \otimes T^{*0,1}M)
\oplus \Sym^2(T^{*0,1}M)$.  The $\OO(2m)$-equivariant maps from this
full space to Hermitian forms are strictly fewer than the
$\UU(m)$-equivariant ones, because $\UU(m)$ identifies the
bidegree $(1,1)$ part as a separate summand admitting its own
equivariant map.
\end{remark}

\begin{definition}[Symbol of a complex rule]
\label{def:symbol-complex}
The \emph{principal symbol} of a complex rule $\Phi$ of order $k$
at a seam~$\varphi$ is the leading-order part of $\Phi$ as a map from
the $(k,0)+(0,k)+(1,k-1)+\cdots$ bidegree parts of the jet to
$\mathcal{H}(M)$.  For a Hermitian-form-valued rule, the
symbol at a cotangent direction $(\xi, \bar\xi) \in T^{*1,0}_xM
\oplus T^{*0,1}_xM$ is a map
\begin{equation}
\sigma_k(\Phi,\varphi,x) \colon
\bigoplus_{p+q=k} (T^{*1,0}_xM)^{\otimes p}
\otimes (T^{*0,1}_xM)^{\otimes q}
\longrightarrow \Herm(T^*_xM).
\end{equation}
The component at bidegree $(1,1)$---encoding the mixed Hessian
direction $\xi \otimes \bar\xi$---governs the ellipticity of the
linearised rule and is the key quantity in the inverse-design
analysis (\S\ref{sec:linearisation}).
\end{definition}

% ================================================================
\section{Classification of Hermitian-Metric-Generating Rules}
\label{sec:classification}
% ================================================================

We now state and prove the central classification result for the
complex setting: the analogue of the Universal Decomposition of the
real-seam paper~\cite{Ronnback2025}, now over $\mathbb{C}$.

Here and below, we use \emph{isotropic} to mean \emph{unitary-isotropic}:
fibrewise $\UU(m)$-equivariant (equivalently, invariant under changes of
$h$-unitary frames).

\subsection{The representation-theoretic setup}

At a point $x \in M$ (working in $h$-unitary normal coordinates),
the complex 2-jet of a seam $\varphi = u + iv$ provides the data
\[
(\varphi_0,\; v,\; \bar w,\; P,\; B,\; \bar{Q})
\;\in\;
\mathbb{C}
\times \mathbb{C}^m
\times \overline{\mathbb{C}^m}
\times \Sym^2(\mathbb{C}^m)
\times M_{m}(\mathbb{C})
\times \overline{\Sym^2(\mathbb{C}^m)},
\]
where $v_\alpha = \partial_\alpha\varphi(x)$,
$\bar w_{\bar\alpha} = \dbar_{\bar\alpha}\varphi(x)$,
$P_{\alpha\beta} = \partial_\alpha\partial_\beta\varphi(x)$, and
$B_{\alpha\bar\beta} = \partial_\alpha\dbar_{\bar\beta}\varphi(x)$.
The group $\UU(m)$ acts by:
$U \cdot v = U v$, $U \cdot P = UPU^T$, $U \cdot B = UB U^\dagger$,
and with appropriate conjugates on the antiholomorphic components.

The target is the space of Hermitian forms:
$\Herm_+(\mathbb{C}^m) \cong \mathbb{R}\cdot I \oplus \mathfrak{su}(m)$
(trace part plus trace-free part) as a real $\UU(m)$-representation.

The \emph{first fundamental theorem of invariant theory for $\UU(m)$}
\cite{Weyl1946, Procesi2007, GoodmanWallach2009} characterises all polynomial
$\UU(m)$-equivariant maps
\[
F \colon \mathbb{C} \times \mathbb{C}^m \times \overline{\mathbb{C}^m} \times M_m(\mathbb{C})
\longrightarrow \Herm(\mathbb{C}^m).
\]
At first derivative order, a complex seam carries two independent holomorphic
gradients $\partial\varphi$ and $\partial\bar\varphi = \overline{\dbar\varphi}$.
Consequently, the unitary-isotropic quadratic gradient contribution is governed
by a $2\times 2$ Hermitian matrix of coefficients acting on the pair
$(\partial\varphi,\partial\bar\varphi)$.  At second derivative order, the mixed
Hessian $B=\partial\dbar\varphi$ contributes its Hermitian and skew-Hermitian
parts (equivalently, the real and imaginary Levi forms).

\subsection{Building blocks}

\begin{proposition}[Equivariant building blocks over $\mathbb{C}$]
\label{prop:building-blocks-c}
Write $v := \partial\varphi(x) \in \mathbb{C}^m$ and
$\tilde v := \partial\bar\varphi(x) = \overline{\dbar\varphi(x)} \in \mathbb{C}^m$.
The space of $\UU(m)$-equivariant maps
from $\mathbb{C} \oplus \mathbb{C}^m \oplus \mathbb{C}^m \oplus M_m(\mathbb{C})$
to $\Herm(\mathbb{C}^m)$ that are linear in $B$ and at most quadratic in
$(v,\tilde v)$ is spanned by the following generators:
\begin{enumerate}[label=(\roman*)]
\item $\Pi_h \colon (\varphi_0, v, B) \mapsto
  \mathrm{Re}(\varphi_0)\cdot I$
  \quad\emph{(scalar-times-identity)};
\item $\Pi_{11} \colon (v,\tilde v) \mapsto v_\alpha\bar{v}_\beta$
  \quad\emph{(holomorphic gradient-square)};
\item $\Pi_{22} \colon (v,\tilde v) \mapsto \tilde v_\alpha\overline{\tilde v}_\beta$
  \quad\emph{(antiholomorphic gradient-square)};
\item $\Pi_{12,+} \colon (v,\tilde v) \mapsto
  \frac{1}{2}\bigl(v_\alpha\overline{\tilde v}_\beta + \tilde v_\alpha\bar v_\beta\bigr)$
  \quad\emph{(mixed gradient, Hermitian part)};
\item $\Pi_{12,-} \colon (v,\tilde v) \mapsto
  \frac{1}{2i}\bigl(v_\alpha\overline{\tilde v}_\beta - \tilde v_\alpha\bar v_\beta\bigr)$
  \quad\emph{(mixed gradient, skew part converted to Hermitian)};
\item $\Pi_{\mathrm{Lev},+} \colon (\varphi_0, v, B) \mapsto
  \frac{1}{2}\bigl(B_{\alpha\bar\beta} + \overline{B_{\beta\bar\alpha}}\bigr)$
  \quad\emph{(Hermitian part of the Levi form)};
\item $\Pi_{\mathrm{Lev},-} \colon (\varphi_0, v, B) \mapsto
  \frac{1}{2i}\bigl(B_{\alpha\bar\beta} - \overline{B_{\beta\bar\alpha}}\bigr)$
  \quad\emph{(skew-Hermitian part converted to Hermitian)}.
\end{enumerate}
Equivalently, the gradient generators (ii)--(v) may be packaged by a
$2\times 2$ Hermitian coefficient matrix $(\beta_{ij}) \in \Herm(\mathbb{C}^2)$,
and generators (vi) and (vii) may be combined into the
single \emph{complex Levi generator}
$\Pi_{\mathrm{Lev}} \colon (\varphi_0, v, B) \mapsto B_{\alpha\bar\beta}$,
with the understanding that we pass to its Hermitian part to land in
$\Herm(\mathbb{C}^m)$.  The combined form is:
$\operatorname{Re}(\gamma\, B_{\alpha\bar\beta})$ for complex $\gamma$.
\end{proposition}

\begin{proof}[Proof sketch]
By Schur's lemma applied to the $\UU(m)$-representation
$\mathbb{C} \oplus \mathbb{C}^m \oplus \mathbb{C}^m \oplus M_m(\mathbb{C})$, a
$\UU(m)$-equivariant linear map to $\Herm(\mathbb{C}^m)$ is
nonzero only on summands that contain $\Herm(\mathbb{C}^m)$ as a
$\UU(m)$-subrepresentation.

\medskip\noindent
\emph{From $\mathbb{C}$ (trivial $\UU(m)$-action, $\UU(1)$-charged).}
The fixed subspace of $\Herm(\mathbb{C}^m)$ under $\UU(m)$ is
$\mathbb{R}\cdot I$ (real multiples of the identity).  The real
part $\mathrm{Re}(\varphi_0)$ maps to $\mathbb{R}\cdot I$, giving
generator~(i).  (The imaginary part $\mathrm{Im}(\varphi_0)$ maps
to $i\cdot \mathbb{R}\cdot I$, but $i I \notin \Herm(\mathbb{C}^m)$;
it contributes to the skew-Hermitian part and is discarded.)

\medskip\noindent
\emph{From the first derivatives (standard $\UU(m)$-representation).}
There is no \emph{linear} $\UU(m)$-equivariant map
$\mathbb{C}^m \to \Herm(\mathbb{C}^m)$: by Schur's lemma, such
a map would identify $\mathbb{C}^m$ with a subrepresentation
of $\Herm(\mathbb{C}^m)$, but the standard representation does
not appear in $\Herm(\mathbb{C}^m) \cong \mathbb{R} \oplus
\mathfrak{su}(m)$ (both summands are self-dual real representations).

For a complex seam, the 1-jet provides two independent holomorphic gradients
$v = \partial\varphi$ and $\tilde v = \partial\bar\varphi$.
Any Hermitian quadratic form in the pair $(v,\tilde v)$ produces a
$\UU(m)$-equivariant Hermitian matrix via outer products
$v_i\overline{v_j}^*$, $i,j\in\{1,2\}$.
Equivalently, the space of unitary-isotropic quadratic gradient generators is
parameterised by $(\beta_{ij}) \in \Herm(\mathbb{C}^2)$, giving the four real
generators (ii)--(v) of Proposition~\ref{prop:building-blocks-c}.

\medskip\noindent
\emph{From $M_m(\mathbb{C})$ (the adjoint$+$its conjugate).}
As a $\UU(m)$-representation,
$M_m(\mathbb{C}) \cong \Herm(\mathbb{C}^m) \oplus
\mathrm{SkewHerm}(\mathbb{C}^m)$
(Hermitian $\oplus$ skew-Hermitian matrices).
By Schur's lemma,
$\Hom_{\UU(m)}(M_m(\mathbb{C}),\, \Herm(\mathbb{C}^m))$
has exactly two generators: the projection onto $\Herm(\mathbb{C}^m)$
(taking the Hermitian part $\frac{1}{2}(B_{\alpha\bar\beta} + \overline{B_{\beta\bar\alpha}})$) and
the skew-Hermitian projection composed with multiplication by $1/i$
to land in $\Herm(\mathbb{C}^m)$
(giving $\frac{1}{2i}(B_{\alpha\bar\beta} - \overline{B_{\beta\bar\alpha}})$).
These are generators~(vi) and~(vii), or equivalently the single complex
generator $B_{\alpha\bar\beta}$ with $\operatorname{Re}(\gamma\cdot)$
applied.
\end{proof}

\begin{corollary}[Dimension count of the building-block space]
\label{cor:dim-building-blocks-c}
Assume $m\ge 2$.  The real vector space of $\UU(m)$-equivariant maps
\[
F\colon \mathbb{C} \oplus \mathbb{C}^m \oplus \mathbb{C}^m \oplus M_m(\mathbb{C}) \to \Herm(\mathbb{C}^m)
\]
that are linear in $B$ and at most quadratic in $(v,\tilde v)$ has real
dimension $7$.  Equivalently, it is parameterised by
\[
\alpha\in\mathbb{R},\qquad (\beta_{ij})\in\Herm(\mathbb{C}^2),\qquad \gamma\in\mathbb{C}.
\]
\end{corollary}

\begin{remark}[Coordinate-free packaging of the gradient family]
\label{rem:coordfree-gradient}
At a point $x$, the pair of holomorphic gradients
$(v,\tilde v)=(\partial\varphi,\partial\bar\varphi)$ may be viewed as a
complex-linear map $V_x\colon \mathbb{C}^2\to\mathbb{C}^m$ by taking $v$ and
$\tilde v$ as its columns in an $h$-unitary frame.  For
$(\beta_{ij})\in\Herm(\mathbb{C}^2)$, the associated quadratic gradient
contribution is then
\[
V_x\,\beta\,V_x^\dagger\ \in\ \Herm(\mathbb{C}^m),
\]
which is independent of the chosen unitary frame.  In particular, if
$\beta\succeq 0$ as a Hermitian $2\times 2$ matrix, then the resulting gradient
term is positive semidefinite at every point.
\end{remark}

\begin{lemma}[Absence of holomorphic and antiholomorphic Hessian generators]
\label{lem:no-hessian}
The holomorphic Hessian $P_{\alpha\beta} \in \Sym^2(T^{*1,0}M)$
and its conjugate $\bar Q$ do \emph{not} contribute generators to
$\Herm(\mathbb{C}^m)$: there is no $\UU(m)$-equivariant linear
map from $\Sym^2(\mathbb{C}^m)$ to $\Herm(\mathbb{C}^m)$, because
these representations do not match (one is complex, the other is real).
Such maps would require a $\mathbb{C}$-antilinear structure absent
from $\Herm(\mathbb{C}^m)$.  In the language of K\"ahler geometry:
$P_{\alpha\beta} = \partial^2_{\alpha\beta}\varphi$ is a $(2,0)$-tensor
that does not contribute to a $(1,1)$-form without further structure.
This absence is the representation-theoretic reason why K\"ahler
potentials enter metric formulas only through $i\partial\dbar\psi$
and not through $\partial^2\psi$ or $\dbar^2\psi$.
\end{lemma}

\subsection{The Universal Decomposition}

We now translate Proposition~\ref{prop:building-blocks-c} from
fibre algebra into a global geometric statement.

In global notation: a convenient generating set is
\begin{enumerate}[label=(\roman*)]
\item $\omegah$ --- the background fundamental form,
\item the \emph{unitary-isotropic gradient family}, generated by
  \[
  i\partial\varphi \wedge \dbar\bar\varphi,\qquad
  i\partial\bar\varphi \wedge \dbar\varphi,\qquad
  \operatorname{Re}\bigl(i\partial\varphi \wedge \dbar\varphi\bigr),\qquad
  \operatorname{Im}\bigl(i\partial\varphi \wedge \dbar\varphi\bigr),
  \]
  i.e. Hermitian quadratic forms built from the pair
  $(\partial\varphi,\partial\bar\varphi)$,
\item $i\partial\dbar\,\mathrm{Re}(\varphi)$ --- the
  \emph{real Levi form} of the seam,
\item $i\partial\dbar\,\mathrm{Im}(\varphi)$ --- the
  \emph{imaginary Levi form} of the seam.
\end{enumerate}
The Levi generators (iii) and (iv) combine elegantly into
$\operatorname{Re}(\gamma\, i\partial\dbar\varphi)$ for complex
$\gamma = \gamma_+ + i\gamma_-$:
\[
\operatorname{Re}\bigl((\gamma_+ + i\gamma_-)\cdot
i\partial\dbar(u + iv)\bigr)
= \gamma_+\, i\partial\dbar u - \gamma_-\, i\partial\dbar v.
\]

\begin{theorem}[Universal Decomposition of order-$\le 2$ isotropic
  complex metric rules]
\label{thm:classification-complex}
Let $(M^{2m}, h, J)$ be a Hermitian manifold with $m \ge 2$.
Every smooth isotropic natural differential operator
$\Phi \colon J^2_\mathbb{C}(M, \mathbb{C}) \to \mathcal{H}(M)$
that depends at most linearly in the mixed Hessian
$\partial\dbar\varphi$ takes the form
\begin{equation}
\label{eq:universal-complex}
\boxed{
\Phi(j^2_\mathbb{C}\varphi;\, h, J)
= \alpha\,\omegah
  + \beta_{11}\, i\partial\varphi \wedge \dbar\bar\varphi
  + \beta_{22}\, i\partial\bar\varphi \wedge \dbar\varphi
  + \operatorname{Re}\!\bigl(\beta_{12}\, i\partial\varphi \wedge \dbar\varphi\bigr)
  + \operatorname{Re}\!\bigl(\gamma\, i\partial\dbar\varphi\bigr),
}
\end{equation}
where $\alpha,\beta_{11},\beta_{22} \in \mathbb{R}$ and
$\beta_{12},\gamma \in \mathbb{C}$ are smooth functions of
$\UU(m)$-invariant scalar jets of order $\le 2$.  One convenient choice of
invariant feature map is:
\begin{equation}
\label{eq:invariants-complex}
\begin{aligned}
\mathcal{I}_\mathbb{C}=\Bigl(
  &\mathrm{Re}(\varphi),\; \mathrm{Im}(\varphi),\; |\partial\varphi|_h^2,\; |\dbar\varphi|_h^2,\\
  &\mathrm{Re}\,\langle\partial\varphi,\partial\bar\varphi\rangle_h,\;
   \mathrm{Im}\,\langle\partial\varphi,\partial\bar\varphi\rangle_h,\\
  &\Delta_h\,\mathrm{Re}(\varphi),\; \Delta_h\,\mathrm{Im}(\varphi),\; |i\partial\dbar\varphi|_h^2
\Bigr),
\end{aligned}
\end{equation}
where $\Delta_h = \tr_{\omegah}(i\partial\dbar)$ is the complex
Laplacian on functions.

The quantity $|i\partial\dbar\varphi|^2_h$ is the squared Frobenius norm of
the complex matrix $B = \partial\dbar\varphi$, namely
$|i\partial\dbar\varphi|^2_h = \tr(B^\dagger B)$ in any $h$-unitary frame.
For complex seams $\varphi$ the mixed Hessian coefficients
$B_{\alpha\bar\beta}=\partial_\alpha\dbar_{\bar\beta}\varphi$ form a general
complex matrix; only its Hermitian part corresponds to
$i\partial\dbar\,\mathrm{Re}(\varphi)$.

The condition $\Phi > 0$ (positive-definite output) imposes
algebraic constraints on $(\alpha, \beta_{11}, \beta_{22}, \beta_{12}, \gamma)$ at each point,
the most basic being $\alpha > 0$.

When nonlinear dependence on the mixed Hessian is allowed,
additional equivariant terms arise: $i\partial\dbar\varphi \cdot
\omegah^{-1} \cdot (i\partial\dbar\varphi)$ (a degree-2 term in
$B$) and $i\partial\varphi \wedge \dbar\bar\varphi \cdot
|i\partial\dbar\varphi|^2_h$ (mixed); the full ring of equivariants
is generated by these together with~\eqref{eq:universal-complex}.
\end{theorem}

\begin{proof}
\emph{Step 1 (Reduction to fibre).}
By the Peetre--Slov\'ak theorem \cite{Peetre1960, Slovak1988},
a smooth natural operator of order $\le 2$ factors through
$J^2_\mathbb{C}(M,\mathbb{C})$.  By $\UU(m)$-equivariance
(axiom~\ref{item:naturality-c}), the value at a point depends
only on the complex 2-jet $j^2_{x,\mathbb{C}}\varphi$ expressed
in $h$-unitary normal coordinates, and is $\UU(m)$-equivariant.
Equivalently, there is a smooth bundle map $F$ such that
$\Phi = F \circ j^2_\mathbb{C}$.
The problem therefore reduces to classifying $\UU(m)$-equivariant
maps
\[
F \colon
\underbrace{\mathbb{C}}_{\varphi_0}
\times \underbrace{\mathbb{C}^m}_{v}
\times \underbrace{\mathbb{C}^m}_{\tilde v = \partial\bar\varphi}
\times \underbrace{M_m(\mathbb{C})}_{B}
\longrightarrow \Herm(\mathbb{C}^m)
\]
that are smooth and at most linear in $B$.

\emph{Step 2 (Decomposition).}
We decompose the source as a $\UU(m)$-representation and identify
equivariant maps to the target $\Herm(\mathbb{C}^m)$.
This is carried out in the proof of
Proposition~\ref{prop:building-blocks-c}, yielding generators
$\Pi_h$, $\Pi_{11}$, $\Pi_{22}$, $\Pi_{12,+}$, $\Pi_{12,-}$,
$\Pi_{\mathrm{Lev},+}$, and $\Pi_{\mathrm{Lev},-}$.

\emph{Step 3 (Coefficients).}
The coefficients may depend smoothly on any choice of
$\UU(m)$-invariant scalar features of the 2-jet.
At first derivative order this includes the norms and inner products of the
pair $(v,\tilde v) = (\partial\varphi,\partial\bar\varphi)$;
at second derivative order it includes trace invariants of
$B = \partial\dbar\varphi$ (and of words in $B$ and $B^\dagger$ if one works
polynomially).  The tuple~\eqref{eq:invariants-complex} is one convenient
low-order choice.

\emph{Step 4 (Completeness).}
Every $\UU(m)$-equivariant map from
$\mathbb{C} \oplus \mathbb{C}^m \oplus \mathbb{C}^m \oplus M_m(\mathbb{C})$ to
$\Herm(\mathbb{C}^m)$ that is linear in $B$ and polynomial of degree
$\le 2$ in $(v,\tilde v)$ is a linear combination (with invariant coefficients)
of the generators in Proposition~\ref{prop:building-blocks-c}.
This follows from the exhaustive $\UU(m)$-representation-theoretic
analysis: no further copies of $\Herm(\mathbb{C}^m)$ appear in
the relevant tensor products.

\emph{Step 5 (Restoring global notation).}
Generator (i) gives $\mathrm{Re}(\varphi_0)\cdot\omegah$; the gradient generators give
$i\partial\varphi \wedge \dbar\bar\varphi$,
$i\partial\bar\varphi \wedge \dbar\varphi$,
$\operatorname{Re}(i\partial\varphi \wedge \dbar\varphi)$ and
$\operatorname{Im}(i\partial\varphi \wedge \dbar\varphi)$; the Levi generators
give
$\frac{1}{2}\bigl(B_{\alpha\bar\beta} + \overline{B_{\beta\bar\alpha}}\bigr)$ and
$\frac{1}{2i}\bigl(B_{\alpha\bar\beta} - \overline{B_{\beta\bar\alpha}}\bigr)$,
which are
$i\partial\dbar\,\mathrm{Re}(\varphi)$ and
$i\partial\dbar\,\mathrm{Im}(\varphi)$ in global notation.
These combine as $\mathrm{Re}(\gamma\, i\partial\dbar\varphi)$.
The coefficient dependence on
$\mathrm{Re}(\varphi_0)$ and $\mathrm{Im}(\varphi_0)$ is absorbed
into the scalar-valued coefficient function $\alpha$.
\end{proof}

\begin{remark}[Holomorphic seams]
\label{rem:holomorphic-seams}
If $\dbar\varphi = 0$ (i.e.\ $\varphi$ is holomorphic), then
$\partial\bar\varphi = 0$ and $\partial\dbar\varphi = 0$, so the universal
decomposition~\eqref{eq:universal-complex} reduces to
\[
\Phi(j^2_\mathbb{C}\varphi;h,J)
= \alpha\,\omegah
  + \beta_{11}\, i\partial\varphi \wedge \dbar\bar\varphi.
\]
In particular, holomorphic seams eliminate both the mixed-gradient terms and
the Levi-form generator.
\end{remark}

\begin{remark}[Complex dimension one]
\label{rem:m1}
The assumption $m\ge 2$ excludes the Riemann surface case.  When $m=1$,
$\Herm(\mathbb{C})$ is $1$-dimensional over $\mathbb{R}$, so every Hermitian
$(1,1)$-form is pointwise a scalar multiple of $\omegah$.  Consequently, any
order-$\le 2$ complex seam rule produces a conformal metric
$\Phi = f\,\omegah$ for some scalar natural operator $f$ depending on the jet
of $\varphi$.
\end{remark}

\begin{remark}[Phase-rotation symmetry]
\label{rem:phase-symmetry}
In applications with a $\UU(1)$ phase symmetry, one may impose invariance under
constant rotations $\varphi\mapsto e^{i\theta}\varphi$.  If, in addition, the
coefficient functions in~\eqref{eq:universal-complex} are required to depend
only on phase-invariant scalar jets (e.g.\ $|\varphi|$, $|\partial\varphi|_h$,
and other $\theta$-independent contractions), then phase invariance forces
$\beta_{12}=0$ and $\gamma=0$, since the mixed term
$i\partial\varphi\wedge\dbar\varphi$ acquires a factor $e^{2i\theta}$ and the
Levi term $\partial\dbar\varphi$ acquires a factor $e^{i\theta}$.
\end{remark}

\begin{corollary}[Phase-invariant subsector]
\label{cor:phase-invariant-subsector}
Under the phase-rotation hypothesis in Remark~\ref{rem:phase-symmetry}, the
maximal phase-invariant sub-family of the universal rule is
\[
\Phi
= \alpha\,\omegah
  + \beta_{11}\, i\partial\varphi \wedge \dbar\bar\varphi
  + \beta_{22}\, i\partial\bar\varphi \wedge \dbar\varphi,
\]
where $\alpha,\beta_{11},\beta_{22}\in\mathbb{R}$.
Equivalently, among the seven real building-block coefficients, phase
invariance eliminates precisely the four parameters corresponding to
$\Pi_{12,\pm}$ and $\Pi_{\mathrm{Lev},\pm}$.
\end{corollary}

\begin{corollary}[Gradient terms obstruct the K\"ahler condition]
\label{cor:gradient-obstructs-kahler}
Assume $d\omegah=0$ (i.e.\ the background Hermitian form is K\"ahler) and that
$\alpha,\beta_{11},\beta_{22}\in\mathbb{R}$ and $\beta_{12},\gamma\in\mathbb{C}$
are constant.
Then the universal rule~\eqref{eq:universal-complex} produces a $d$-closed
$(1,1)$-form for every seam $\varphi$ if and only if
$\beta_{11}=\beta_{22}=0$ and $\beta_{12}=0$.
In particular, the Levi-form term $\mathrm{Re}(\gamma\,i\partial\dbar\varphi)$
is automatically K\"ahler, while the gradient-tensor generators are the
precise obstruction.
\end{corollary}

\begin{proof}
Since $d=\partial+\dbar$ and $\partial^2=\dbar^2=0$, we have
$d\,\mathrm{Re}(\gamma\,i\partial\dbar\varphi)=0$ for every $\varphi$ and every
constant $\gamma$.
By contrast,
\[
d\bigl(i\partial\varphi\wedge\dbar\bar\varphi\bigr)
= -i\partial\varphi\wedge\partial\dbar\bar\varphi
  + i\dbar\partial\varphi\wedge\dbar\bar\varphi,
\]
which is generically nonzero (and similarly for the other gradient-tensor
generators).  Hence $d\Phi=0$ for all seams forces
$\beta_{11}=\beta_{22}=\beta_{12}=0$.
The converse is immediate when $d\omegah=0$.
\end{proof}

\begin{corollary}[Forced factorisation through invariants]
\label{cor:forced-arch-c}
Under the hypotheses of Theorem~\ref{thm:classification-complex},
the value of $\Phi(j^2_\mathbb{C}\varphi;\, h,J)$ at $x \in M$
depends on the 2-jet of $\varphi$ at $x$ only through
$\UU(m)$-invariant scalar functions of that jet and the basis forms
$\omegah(x)$,
$(i\partial\varphi \wedge \dbar\bar\varphi)(x)$,
$(i\partial\bar\varphi \wedge \dbar\varphi)(x)$,
$\operatorname{Re}(i\partial\varphi \wedge \dbar\varphi)(x)$,
$\operatorname{Im}(i\partial\varphi \wedge \dbar\varphi)(x)$,
$(i\partial\dbar\,\mathrm{Re}(\varphi))(x)$, and
$(i\partial\dbar\,\mathrm{Im}(\varphi))(x)$.

Equivalently, any parametric family of such rules (e.g., a neural
architecture constrained to represent a smooth isotropic natural
complex operator of order $\le 2$ linear in $\partial\dbar\varphi$)
takes the canonical form
\[
\begin{aligned}
\Phi_\theta(j^2_\mathbb{C}\varphi;\, h,J)
={}&\alpha_\theta(\mathcal{I}_\mathbb{C})\,\omegah
+ \beta_{11,\theta}(\mathcal{I}_\mathbb{C})\,i\partial\varphi \wedge \dbar\bar\varphi \\
&+ \beta_{22,\theta}(\mathcal{I}_\mathbb{C})\,i\partial\bar\varphi \wedge \dbar\varphi \\
&+ \operatorname{Re}\!\bigl(\beta_{12,\theta}(\mathcal{I}_\mathbb{C})\,i\partial\varphi \wedge \dbar\varphi\bigr) \\
&+ \operatorname{Re}\!\bigl(\gamma_\theta(\mathcal{I}_\mathbb{C})\,i\partial\dbar\varphi\bigr)
\end{aligned}
\]
for smooth coefficient functions
$(\alpha_\theta, \beta_{11,\theta}, \beta_{22,\theta})\colon \mathbb{R}^N \to \mathbb{R}^3$
and $(\beta_{12,\theta}, \gamma_\theta)\colon \mathbb{R}^N \to \mathbb{C}^2$,
for some choice of invariant feature map $\mathcal{I}_\mathbb{C}\colon J^2_\mathbb{C}\to\mathbb{R}^N$.
\end{corollary}

\subsection{Recovery of known structures}

Every Hermitian metric-generating construction in the literature
is a specialisation of~\eqref{eq:universal-complex}:

\begin{example}[K\"ahler metric from a K\"ahler potential]
\label{ex:kahler}
$\alpha = 0$, $\beta_{11} = \beta_{22} = 0$, $\beta_{12} = 0$, $\gamma = 1$ (real), $\varphi = \psi$
real (plurisubharmonic).
Output: $\Phi = i\partial\dbar\psi$.
This is the fundamental $(1,1)$-form of the K\"ahler metric
associated to the potential $\psi$.  Strict plurisubharmonicity of
$\psi$ ($i\partial\dbar\psi > 0$) is precisely the positivity
condition for $\Phi$.
\end{example}

\begin{proposition}[Chern--Ricci form of a Levi-form metric]
\label{prop:chern-ricci-levi}
Let $\omega = i\partial\dbar\psi$ be a K\"ahler metric written in local
holomorphic coordinates as
$\omega = i\, g_{\alpha\bar\beta}\, dz^\alpha\wedge d\bar z^\beta$ with
$g_{\alpha\bar\beta} = \partial_\alpha\dbar_{\bar\beta}\psi$.
Then its Chern--Ricci form is
\[
\Ric(\omega) = -i\partial\dbar\log\det(g_{\alpha\bar\beta}).
\]
\end{proposition}

\begin{remark}[Calabi--Yau as a seam equation]
Fix a compact K\"ahler manifold $(M,\omega_0,J)$.  Writing a metric in the same
K\"ahler class as $\omega = \omega_0 + i\partial\dbar\psi$ turns the problem of
prescribing the volume form (equivalently, the Ricci form within the class)
into a PDE for the scalar seam $\psi$.  Concretely, prescribing
$\omega^m = e^F\,\omega_0^m$ yields the complex Monge--Amp\`ere equation
\[
(\omega_0 + i\partial\dbar\psi)^m = e^F\,\omega_0^m,
\]
with $\psi$ determined up to an additive constant.  The Calabi--Yau theorem
asserts existence and uniqueness (up to constant) under the usual compatibility
condition on total volume \cite{Yau1978}.
\end{remark}

\begin{example}[Fubini--Study metric]
\label{ex:fs}
Let $M = \mathbb{C}^m$ with background
$h = \delta$ (flat), and take the seam
$\varphi = \log(1 + |z|^2)$ (real, plurisubharmonic).
Then $\alpha = 0$, $\beta_{11} = \beta_{22} = 0$, $\beta_{12} = 0$, $\gamma = 1$, and
$\Phi = i\partial\dbar\log(1+|z|^2)$
is the Fubini--Study metric on $\mathbb{CP}^m$ in affine
coordinates.
\end{example}

\begin{example}[Hermitian conformal rule]
\label{ex:herm-conformal}
$\alpha = e^{2\psi}$ (for real seam $\psi$), $\beta_{11} = \beta_{22} = 0$, $\beta_{12} = 0$, $\gamma = 0$.
Output: $\Phi = e^{2\psi}\omegah$.
This is a conformal rescaling of the background Hermitian metric:
it preserves the complex structure $J$ and the conformal class
$[\omegah]$ but changes the pointwise scale.
\end{example}

\begin{example}[Holomorphic gradient metric]
\label{ex:grad-metric}
$\alpha = 0$, $\beta_{11} = 1$, $\beta_{22} = 0$, $\beta_{12} = 0$, $\gamma = 0$.
Output: $\Phi = i\partial\varphi \wedge \dbar\bar\varphi$.
This is a positive semidefinite Hermitian form of rank $\le 1$,
vanishing precisely at the points where $\partial\varphi = 0$
(the \emph{holomorphic critical points} of $\varphi$).
For holomorphic $\varphi$, these coincide with the algebraic
critical points; for a general complex seam, $\partial\varphi = 0$
and $\dbar\varphi = 0$ may be independent conditions.
\end{example}

\begin{example}[Fisher--Rao metric of a complex exponential family]
\label{ex:fisher-complex}
Let $\{\,p(x|\theta)\,\}$ be an exponential family parametrised by
complex $\theta \in \Theta \subset \mathbb{C}^m$, with log-partition
function $A(\theta) \in \mathbb{R}$ (a real function of complex
parameters).  The complex Fisher information metric on $\Theta$ is
\[
g^{\mathrm{Fisher}}_{\alpha\bar\beta}(\theta)
= \partial_\alpha\dbar_{\bar\beta}\,A(\theta),
\]
which is precisely the Levi-form rule ($\alpha = 0$, $\beta_{11} = \beta_{22} = 0$, $\beta_{12} = 0$,
$\gamma = 1$) applied to the real seam $\varphi = A(\theta)$ on the
flat background $(\Theta, \delta)$.  Strict convexity of $A$ (a
standard property of exponential families) is equivalent to
strict plurisubharmonicity of $A$ viewed as a function of complex
parameters, and hence to positive-definiteness of the resulting
Hermitian metric.
\end{example}

\begin{example}[Mixed Hermitian-Levi rule]
\label{ex:mixed-c}
$\alpha = e^{|\varphi|^2}$, $\beta_{11} = \beta_{22} = 0$, $\beta_{12} = 0$, $\gamma = \lambda \in \mathbb{R}$.
Output: $\Phi = e^{|\varphi|^2}\omegah + \lambda\, i\partial\dbar
\mathrm{Re}(\varphi)$.
The background term ensures positive-definiteness for small $|\lambda|$,
while the Levi-form term introduces curvature.  This is the complex
analogue of the mixed conformal--Hessian rule of~\cite{Ronnback2025},
and arises naturally in the study of curvature perturbations of
K\"ahler metrics.
\end{example}

\begin{corollary}[No globally K\"ahler pure Levi metrics on
  closed manifolds with non-K\"ahler fundamental class]
\label{cor:no-global-levi}
Let $(M^{2m},h,J)$ be a closed Hermitian manifold, and let
$\varphi \in C^\infty(M,\mathbb{C})$.  If the de Rham cohomology
class $[\omegah]$ is not in the image of $\partial\dbar\colon
C^\infty(M,\mathbb{C}) \to \mathcal{H}(M)$---i.e.\ if $\omegah$
is not globally $\partial\dbar$-exact---then the pure Levi-form rule
$g = i\partial\dbar\,\mathrm{Re}(\varphi)$ cannot produce a metric
cohomologous to $\omegah$.
\end{corollary}

\begin{proof}
Since $i\partial\dbar\,\mathrm{Re}(\varphi)$ is globally
$\partial\dbar$-exact (equal to $\frac{i}{2}(dd^c)\mathrm{Re}(\varphi)$
where $d^c = i(\dbar - \partial)$), it represents the zero class in
de Rham cohomology.  A metric cohomologous to $\omegah$ would require
$[\omegah] = 0$ in $H^{1,1}(M,\mathbb{R})$, which by assumption does
not hold.
\end{proof}

This is the complex analogue of the obstruction for the pure
Hessian rule on closed real manifolds \cite{Ronnback2025}.

\subsection{The degeneration locus for complex rules}

The degeneration locus inherits a richer structure in the complex
setting, reflecting the bidegree decomposition of the jet.

\begin{definition}[Complex degeneration locus]
\label{def:degen-c}
For a complex rule $\Phi$ and a seam $\varphi$, the
\emph{complex degeneration locus} is
\[
\mathcal{D}_\mathbb{C}(\Phi,\varphi)
:= \bigl\{\, x \in M :
  \text{the fibre map }
  D\Phi_\varphi(x) \colon J^k_{x,\mathbb{C}} \to
  \mathcal{H}^+(T^*_xM)
  \text{ is not surjective}\bigr\}.
\]
For some basic rules:
\begin{enumerate}[label=(\arabic*)]
\item \textbf{Hermitian conformal rule} $(\alpha = e^{2\psi},
  \beta_{11} = \beta_{22} = \beta_{12} = \gamma = 0)$: $\mathcal{D}_\mathbb{C} = \emptyset$
  (the conformal factor $e^{2\psi}$ is everywhere positive).
\item \textbf{Holomorphic gradient rule} $(\beta_{11} = 1,
  \alpha = \beta_{22} = \beta_{12} = \gamma = 0)$: $\mathcal{D}_\mathbb{C} =
  \{\partial\varphi = 0\}$ (the holomorphic critical points),
  which for a holomorphic $\varphi$ coincides with the
  algebraic zero set $\{\varphi' = 0\}$.
\item \textbf{Pure Levi-form rule} $(\gamma = 1,
  \alpha = \beta_{11} = \beta_{22} = \beta_{12} = 0)$: $\mathcal{D}_\mathbb{C} =
  \{\det_h(i\partial\dbar\,\mathrm{Re}(\varphi)) = 0\}$, the
  \emph{Levi-degenerate set} of the seam.
  For a plurisubharmonic seam, this is the boundary of the
  strictly plurisubharmonic domain.
\end{enumerate}
\end{definition}

\begin{remark}[Kernel distribution for the holomorphic gradient form]
Away from $\{\partial\varphi = 0\}$, the form
$i\partial\varphi\wedge\dbar\bar\varphi$ has rank $1$ and its kernel is the
complex hyperplane distribution $\ker(\partial\varphi)\subset T^{1,0}M$.
If $\varphi$ is holomorphic and a submersion, this distribution integrates
to the holomorphic foliation by level hypersurfaces of $\varphi$.
\end{remark}

\begin{remark}[Holomorphic critical points vs.\ real critical points]
The degeneration locus of the holomorphic gradient rule is
$\{\partial\varphi = 0\}$, which is genuinely different from
the real critical set $\{\nabla_{\mathbb{R}}\varphi = 0\}$:
since $\nabla_{\mathbb{R}}\varphi = 0$ if and only if both
$\partial\varphi = 0$ and $\dbar\varphi = 0$, we have
$\mathcal{D}_\mathbb{C} \supseteq \Crit_\mathbb{R}(\varphi)$
with equality only when $\varphi$ is holomorphic.  For a
general complex seam, $\partial\varphi = 0$ does not require
$\dbar\varphi = 0$, so the complex degeneration locus may be
strictly larger than the real critical set.  This enlargement
is precisely the additional geometric information the complex
structure detects.
\end{remark}

% ================================================================
\section{The Linearisation of a Complex Rule}
\label{sec:linearisation}
% ================================================================

Given a complex rule $\Phi$ and a background seam $\varphi_0$,
the first-order response to a perturbation $\dot\varphi =
\dot u + i\dot v$ is
\begin{equation}
\label{eq:linearisation-c}
D\Phi_{\varphi_0}(\dot\varphi) :=
\left.\frac{d}{dt}\right|_{t=0}
\Phi(\varphi_0 + t\dot\varphi;\, h, J).
\end{equation}
This is a $\mathbb{R}$-linear differential operator
$D\Phi_{\varphi_0}\colon C^\infty(M,\mathbb{C}) \to
\Gamma(\mathcal{H}(M))$ (into the bundle of Hermitian forms,
not necessarily positive-definite).

\begin{proposition}[Linearisation of the universal decomposition]
\label{prop:lin-c}
For the rule
\[
\Phi = \alpha\,\omegah
  + \beta_{11}\, i\partial\varphi \wedge \dbar\bar\varphi
  + \beta_{22}\, i\partial\bar\varphi \wedge \dbar\varphi
  + \operatorname{Re}\!\bigl(\beta_{12}\, i\partial\varphi \wedge \dbar\varphi\bigr)
  + \operatorname{Re}\!\bigl(\gamma\,i\partial\dbar\varphi\bigr),
\]
with $\alpha,\beta_{11},\beta_{22} \in \mathbb{R}$ and
$\beta_{12},\gamma \in \mathbb{C}$ constant, the linearisation at any
background seam $\varphi_0$ is:
\[
\begin{aligned}
D\Phi_{\varphi_0}(\dot\varphi)
={}&\beta_{11}\, i\bigl(\partial\dot\varphi \wedge \dbar\bar\varphi_0
  + \partial\varphi_0 \wedge \dbar\overline{\dot\varphi}\bigr) \\
&+ \beta_{22}\, i\bigl(\partial\overline{\dot\varphi} \wedge \dbar\varphi_0
  + \partial\bar\varphi_0 \wedge \dbar\dot\varphi\bigr) \\
&+ \operatorname{Re}\!\bigl(\beta_{12}\, i(\partial\dot\varphi \wedge \dbar\varphi_0
  + \partial\varphi_0 \wedge \dbar\dot\varphi)\bigr) \\
&+ \operatorname{Re}\bigl(\gamma\, i\partial\dbar\dot\varphi\bigr)
 + \text{(coefficient variation terms)}.
\end{aligned}
\]
In particular:
\begin{enumerate}[label=(\alph*)]
\item \textbf{Conformal rule} $(\beta_{11} = \beta_{22} = \beta_{12} = \gamma = 0)$:
  $D\Phi_{\varphi_0}(\dot\varphi) = (\partial_u\alpha)\dot u\,\omegah
  + (\partial_v\alpha)\dot v\,\omegah$, a zeroth-order operator.
\item \textbf{Holomorphic gradient rule} $(\alpha = \beta_{22} = \beta_{12} = \gamma = 0,\, \beta_{11} = 1)$:
  at $\varphi_0$ with $\partial\varphi_0 \neq 0$,
  $D\Phi_{\varphi_0}(\dot\varphi) = i(\partial\dot\varphi \wedge
  \dbar\bar\varphi_0 + \partial\varphi_0 \wedge \dbar\overline{\dot\varphi})$,
  a first-order operator.  At $\partial\varphi_0 = 0$, the
  linearisation vanishes: the rule degenerates at holomorphic
  critical points.
\item \textbf{Pure Levi-form rule} $(\alpha = \beta_{11} = \beta_{22} = \beta_{12} = 0,\, \gamma = 1)$:
  $D\Phi_{\varphi_0}(\dot\varphi) = i\partial\dbar\dot u$,
  a second-order overdetermined-elliptic operator.  Its
  kernel consists of functions $\dot\varphi$ whose real part
  is pluriharmonic: $\ker D\Phi|_{\gamma=1} \supseteq
  \{\dot\varphi \in C^\infty(M,\mathbb{C}) :
  \partial\dbar\,\mathrm{Re}(\dot\varphi) = 0\}$.
\end{enumerate}
\end{proposition}

\begin{corollary}[Kernel rigidity for the Levi linearisation]
\label{cor:levi-kernel-rigidity}
Let $(M,\omega_0,J)$ be a compact K\"ahler manifold and consider the pure
Levi-form rule $\Phi(\varphi)=i\partial\dbar\,\mathrm{Re}(\varphi)$.  For
real-valued variations $\dot\varphi=\dot u\in C^\infty(M,\mathbb{R})$,
\[
\ker(D\Phi_{\varphi_0})\cap C^\infty(M,\mathbb{R}) = \{\text{constants}\}.
\]
\end{corollary}

\begin{proof}
If $D\Phi_{\varphi_0}(\dot u)=i\partial\dbar\dot u=0$, then $\dot u$ is
pluriharmonic.  Tracing with respect to $\omega_0$ gives
$\Delta_{\omega_0}\dot u=0$, hence $\dot u$ is constant by the maximum
principle.
\end{proof}

\begin{remark}[The Levi-form generator and ellipticity]
\label{rem:ellipticity}
For the Levi-form term $\operatorname{Re}(\gamma\,i\partial\dbar\varphi)$, the
linearisation contributes $\dot\varphi \mapsto \operatorname{Re}(\gamma\,i\partial\dbar\dot\varphi)$.
The principal symbol at a covector
$\xi \in T^{*1,0}_xM$ is:
\[
\sigma\bigl(\operatorname{Re}(\gamma\,i\partial\dbar\cdot)\bigr)(\xi)\cdot\dot\varphi
= \xi \otimes \bar\xi \cdot \mathrm{Re}(\gamma\,\dot\varphi),
\]
the $(1,1)$-form $\xi \wedge \bar\xi$ (positive semidefinite, rank 1).
In particular, restricting to real variations $\dot\varphi=\dot u$ gives a
symbol proportional to $\mathrm{Re}(\gamma)\,\xi\otimes\bar\xi$, so the Levi
term degenerates at the level of the principal symbol precisely when
$\mathrm{Re}(\gamma)=0$.
This is the complex analogue of the real symbol $\xi \otimes \xi$ for
the Hessian generator of~\cite{Ronnback2025}, with $\xi \in T^*_xM$
replaced by the $(1,0)$-component $\xi \in T^{*1,0}_xM$.  The key
structural difference is that the complex symbol is Hermitian
(satisfying $(\xi \otimes \bar\xi)^\dagger = \bar\xi \otimes \xi =
\xi \otimes \bar\xi$ up to the Hermitian identification), while the
real symbol is only symmetric.  This Hermitian positivity of the
symbol implies that, as an operator on \emph{real} variations
$\dot\varphi=\dot u$, the Levi linearisation is overdetermined-elliptic: the
symbol $\xi\otimes\bar\xi\cdot \dot u$ is injective as a map
$\mathbb{R}\to\Herm(T^{1,0}_xM)$ whenever $\xi\neq 0$.
The term \emph{subelliptic} is used in a different, genuinely analytic sense
(e.g.\ Kohn--Nirenberg estimates) and does not follow from this symbol
computation alone.  On K\"ahler backgrounds, additional Hodge-theoretic
structure (including K\"ahler identities such as $\Box_{\partial}=\Box_{\dbar}$)
can lead to sharper spectral estimates, but this is distinct from
subellipticity.
\end{remark}

% ================================================================
\section{Discussion: Towards a Universality Theory}
\label{sec:universality}
% ================================================================

The universality question for complex seams is: when can scalar
seam data $\varphi$ generate, via rule~\eqref{eq:universal-complex},
an open neighbourhood of all Hermitian metrics on $M$ (up to
biholomorphism)?

The situation is markedly different from the real case.  On a compact
K\"ahler manifold, the $\partial\dbar$-lemma (a consequence of Hodge
theory) asserts that every $\partial\dbar$-closed real $(1,1)$-form
in a fixed de Rham class is $\partial\dbar$-exact, meaning
$\omega = \omega_0 + i\partial\dbar\psi$ for a real potential $\psi$.

\begin{theorem}[Local universality in a K\"ahler class]
\label{thm:local-universality-kahler}
Let $(M,\omega_0,J)$ be a compact K\"ahler manifold.  If $\omega$ is any
K\"ahler metric with $[\omega]=[\omega_0]$ in de Rham cohomology, then
\[
\omega = \omega_0 + i\partial\dbar\psi
\]
for some $\psi\in C^\infty(M,\mathbb{R})$, unique up to addition of a constant.
\end{theorem}

\begin{proof}[Proof sketch]
Since $[\omega]=[\omega_0]$, the real $(1,1)$-form $\eta:=\omega-\omega_0$ is
exact, hence in particular $d$-closed.  On a compact K\"ahler manifold the
$\partial\dbar$-lemma implies that any exact real $(1,1)$-form is
$\partial\dbar$-exact, so $\eta=i\partial\dbar\psi$ for a real potential
$\psi$.  If also $\omega=\omega_0+i\partial\dbar\psi'$, then
$i\partial\dbar(\psi-\psi')=0$, so tracing with respect to $\omega_0$ gives
$\Delta_{\omega_0}(\psi-\psi')=0$ and thus $\psi-\psi'$ is constant by the
maximum principle.
\end{proof}

This is a \emph{local universality} statement for the Levi-form rule within a
fixed K\"ahler class.  It should be contrasted with the real-seam case, where
universal decomposition at order 2 falls far short of reaching all metrics.


For the full complex seam $\varphi = u + iv$ and the complete
rule~\eqref{eq:universal-complex}, the universality question
acquires a K\"ahler-class component (governed by $i\partial\dbar u$)
and a Hermitian-but-non-K\"ahler component (governed by
$i\partial\dbar v$ and the gradient-tensor term).  We state the
following as open problems, motivated by the real-seam analogues
proved in~\cite{Ronnback2025}:

\begin{openproblem}[Complex local universality]
On a compact K\"ahler manifold $(M,\omega_0, J)$, is the extended
complex seam map
\[
F(\varphi,\psi_\phi)
= (\mathrm{biholomorphism})^*\,\bigl[\,
  \alpha\,\omegah
  + \beta_{11}\,i\partial\varphi\wedge\dbar\bar\varphi
  + \beta_{22}\,i\partial\bar\varphi\wedge\dbar\varphi
  + \operatorname{Re}\!\bigl(\beta_{12}\,i\partial\varphi\wedge\dbar\varphi\bigr)
  + \operatorname{Re}(\gamma\,i\partial\dbar\varphi)
  \,\bigr]
\]
locally surjective onto the cone of Hermitian metrics near any
non-degenerate background?  In particular, does the gradient-tensor
family $(\beta_{ij})\in\Herm(\mathbb{C}^2)$ play the same
role of navigating the Teichm\"uller directions in the complex
setting as it does for the torus in the real setting?
\end{openproblem}

\begin{openproblem}[Complex vs.\ real universality]
Is the complex seam map~\eqref{eq:universal-complex}
``more universal'' than the real seam map in the sense that its image
has smaller codimension in the space of Hermitian metrics?
The $\partial\dbar$-lemma suggests that the Levi-form generator alone
covers the entire K\"ahler class---an infinite-dimensional subspace
of all Hermitian metrics---which is a stronger universality result
than any of the real-seam generators achieve at order $\le 2$.
Make this precise.
\end{openproblem}

% ================================================================
\section{Conclusion}
% ================================================================

The promotion of the seam from a real-valued to a complex-valued
function, combined with the tightening of the naturality group from
$\OO(2m)$ to $\UU(m)$, produces one genuinely new generator in the
classification of isotropic Hermitian-metric rules at order $\le 2$:
the complex Levi form $\operatorname{Re}(\gamma\, i\partial\dbar\varphi)$
for complex $\gamma$.  This single additional generator subsumes the
entire theory of K\"ahler metrics from potentials, the Fisher--Rao
metric of complex exponential families, and the Fubini--Study geometry
of quantum state spaces.

The framework developed here retains all three structural properties
of the real case---classification, Fredholm deformation theory, and
a well-posed universality question---while exploiting the tighter
$\UU(m)$-equivariance to access geometric structures (holomorphic
gradient forms, Levi degeneracy sets, $\partial\dbar$-exact deformations)
that are invisible to the real $\OO(n)$ theory.

\backmatter

\bmhead{Acknowledgements}

Not applicable.

\section*{Declarations}

\bmhead{Funding}

Not applicable.

\bmhead{Competing interests}

The author declares no competing interests.

\bmhead{Use of large language models (LLMs)}

During preparation of this manuscript, the author used large language
models (Claude Opus~4.6, GPT-5.2, Grok~4.20, and Gemini~3.1 Pro) to
assist with drafting and editing text and with exploring alternative
expositions. The author reviewed and edited the resulting content and
takes full responsibility for the manuscript.

\bmhead{Ethics approval and consent to participate}

Not applicable.

\bmhead{Consent for publication}

Not applicable.

\bmhead{Data availability}

Not applicable.

\bmhead{Materials availability}

Not applicable.

\bmhead{Code availability}

Not applicable.

\bmhead{Author contribution}

The author wrote the manuscript.

\begin{thebibliography}{99}

\bibitem{Ronnback2025}
L.~R\"onnb\"ack,
``Seam Geometry: A Jet-Theoretic Framework for the Scalar Generation
of Geometric Structures,''
\textit{preprint}, 2025.

\bibitem{Peetre1960}
J.~Peetre,
``Rectification \`a l'article `Une caract\'erisation abstraite des
op\'erateurs diff\'erentiels'\,''
\textit{Math.\ Scand.} \textbf{8}, 116--120 (1960).

\bibitem{Slovak1988}
J.~Slov\'ak,
``Peetre theorem for nonlinear operators,''
\textit{Ann.\ Global Anal.\ Geom.} \textbf{6}(3), 273--283 (1988).

\bibitem{Weyl1946}
H.~Weyl,
\textit{The Classical Groups: Their Invariants and Representations},
Princeton University Press, 1946.

\bibitem{Procesi2007}
C.~Procesi,
\textit{Lie Groups: An Approach through Invariants and
Representations},
Springer, 2007.

\bibitem{GoodmanWallach2009}
R.~Goodman and N.~R.~Wallach,
\textit{Symmetry, Representations, and Invariants},
Graduate Texts in Mathematics~\textbf{255},
Springer, 2009.

\bibitem{Yau1978}
S.-T.~Yau,
``On the Ricci curvature of a compact K\"ahler manifold and the complex
Monge--Amp\`ere equation, I,''
\textit{Comm.\ Pure Appl.\ Math.} \textbf{31}(3), 339--411 (1978).

\end{thebibliography}

\end{document}
